% ===============================================================
% Section I  --  Introduction and Problem Formulation
% Owner: Integration and Reporting Lead
% ===============================================================
\section{Introduction and Problem Formulation}
\label{sec:intro}

Recent progress in natural language processing has been driven largely
by very large foundation models. In practice, however, many
organisations cannot afford, govern, or justify hosting such models for
narrow, bounded tasks. This has motivated growing interest in
\emph{small language models} (SLMs): compact, deployable models that can
be adapted to a specific domain under realistic constraints on cost,
privacy, latency, and maintainability. The central question this project
addresses is therefore not how to build the largest possible system, but
how to make a small, constrained model reliable and useful on a
well-defined domain task.

\subsection{The Task}
We build a Shakespeare-aware language system that lets a user with little
or no prior exposure to Shakespeare interrogate three plays,
\emph{Hamlet}, \emph{Macbeth}, and \emph{Romeo and Juliet}, and receive
explanations that are accurate, beginner-friendly, and traceable to the
source text. Concretely, the system must support four capabilities:
explain characters, events, and themes; answer contextual questions
about motivations and events; retrieve and display the source passages
that support each answer; and produce short, clearly-labelled responses
in a Shakespearean style without presenting them as fact.

\subsection{Why This Is a Domain Adaptation Problem}
Shakespearean English departs sharply from contemporary English in
vocabulary, syntax, idiom, and cultural reference. A pretrained model
that is fluent in modern English does not, by default, map a plain
modern question onto the archaic language that actually answers it. As
the course material on language modelling emphasises, words carry meaning
through the contexts in which they appear, and a useful system must be
able to retrieve relevant information \emph{even when the query and the
source do not share the same words}. Bridging that gap, between a
beginner's modern phrasing and the early-modern text, is precisely the
domain adaptation challenge here. We address it through retrieval and
prompt design rather than by retraining the model, so that the model's
answers are anchored to evidence drawn from the plays themselves.

\subsection{Intended User and Design Implications}
The intended user is a newcomer who has not studied Shakespeare. This
single assumption shapes the entire design. Explanations must avoid
assuming literary background; every substantive answer must be backed by
visible evidence so the user can verify it; and creative, stylised output
must be unmistakably separated from factual explanation so a beginner is
never misled into treating an invented line as a quotation.

\subsection{Constraints}
Three constraints framed our design decisions. \emph{Compute:} the system
must run on a standard laptop without specialised hardware, which ruled
out large models and steered us toward a compact, locally-hosted model
and cached indices. \emph{Architecture:} a retrieval-augmented generation
pipeline is mandatory, and retrieved passages must demonstrably inform
generation rather than decorate it. \emph{Scope:} pretraining or
full-scale fine-tuning is out of scope; adaptation is achieved through
retrieval and prompting. Working within these constraints, the project's
goal is a system that is well-designed, well-evaluated, and clearly
explained, rather than merely large or complex.

\subsection{Contributions and Report Structure}
The remainder of this report is organised around the system we built.
Section~\ref{sec:dataset} describes how the corpus was prepared, chunked,
and enriched. Section~\ref{sec:methodology} sets out the baseline and
RAG methodology and justifies the choice of model.
Section~\ref{sec:system} describes the end-to-end system, its command-line
interface, and how to run it. Sections~\ref{sec:evaluation}
and~\ref{sec:results} present the evaluation protocol, results, and
failure analysis. Section~\ref{sec:genai} documents our use of generative
AI, and Section~\ref{sec:conclusion} concludes.
